% Problem 4 — Subharmonicity of 1/Phi_n under finite free convolution
% Solution by Claude (Opus 4.7), 2026-04-30
\documentclass[12pt]{article}
\usepackage{fullpage,amsmath,amssymb,amsthm,hyperref}
\newtheorem{theorem}{Theorem}
\newtheorem{proposition}[theorem]{Proposition}
\newtheorem{lemma}[theorem]{Lemma}
\theoremstyle{definition}
\newtheorem{definition}[theorem]{Definition}
\newtheorem{remark}[theorem]{Remark}

\title{Solution to Problem 4: Subharmonicity of $1/\Phi_n$ under Finite Free Convolution}
\author{}
\date{}

\begin{document}
\maketitle

\section*{Statement and answer}

For a monic real-rooted polynomial $p(x)=\prod_{i=1}^n(x-\lambda_i)$ of degree $n$ with distinct roots, the \emph{finite free Fisher information} is
\[
\Phi_n(p) \;:=\; \sum_{i=1}^n\Bigl(\sum_{j\neq i}\frac{1}{\lambda_i-\lambda_j}\Bigr)^2,
\]
and $\Phi_n(p):=+\infty$ if $p$ has a multiple root. For monic polynomials of degree $n$,
\[
p(x)=\sum_{k=0}^n a_kx^{n-k},\qquad q(x)=\sum_{k=0}^n b_kx^{n-k},\qquad a_0=b_0=1,
\]
the \emph{finite free convolution} $p\boxplus_n q$ is the monic degree-$n$ polynomial whose coefficients $c_k$ are
\begin{equation}\label{eq:ck}
c_k \;=\; \sum_{i+j=k}\frac{(n-i)!(n-j)!}{n!\,(n-k)!}\,a_ib_j,\qquad k=0,1,\dots,n.
\end{equation}
The question asks whether
\begin{equation}\label{eq:main}
\frac{1}{\Phi_n(p\boxplus_n q)} \;\geq\; \frac{1}{\Phi_n(p)} \;+\; \frac{1}{\Phi_n(q)}
\end{equation}
holds for all monic real-rooted $p,q$ of degree $n$.

\smallskip
\noindent\textbf{Answer.} We prove the inequality \eqref{eq:main} in full for $n=1$ and $n=2$, with equality in both cases. We provide complete proofs of the underlying score-formula lemma (Part~1) and the explicit $n=2$ computation (Part~2). The general case $n\geq 3$ is the finite-free analogue of Voiculescu's superadditivity of the inverse free Fisher information; we explain in Part~3 precisely what additional ingredients are required to extend our proof to that range, and provide the strongest partial result we can derive in full.

\section{Part 1: The score formula}

In this section we prove a clean closed-form expression for the inner sum that defines the finite free score, in terms of the polynomial $p$ and its derivatives.

\begin{lemma}[Score formula]\label{lem:score}
Let $p(x)=\prod_{i=1}^n(x-\lambda_i)$ be a monic polynomial of degree $n\geq 2$ with simple (pairwise distinct) real roots $\lambda_1,\dots,\lambda_n$. Then for every $k\in\{1,\dots,n\}$,
\begin{equation}\label{eq:scoreformula}
\sum_{j\neq k}\frac{1}{\lambda_k-\lambda_j} \;=\; \frac{p''(\lambda_k)}{2\,p'(\lambda_k)}.
\end{equation}
Consequently,
\begin{equation}\label{eq:phiformula}
\Phi_n(p) \;=\; \frac{1}{4}\sum_{k=1}^n\Bigl(\frac{p''(\lambda_k)}{p'(\lambda_k)}\Bigr)^2.
\end{equation}
\end{lemma}

\begin{proof}
Fix $k$ and define
\[
r_k(x) \;:=\; \prod_{j\neq k}(x-\lambda_j).
\]
Then $r_k$ is a monic polynomial of degree $n-1$ and
\begin{equation}\label{eq:factor}
p(x) \;=\; (x-\lambda_k)\,r_k(x).
\end{equation}

\textbf{Step 1: $p'(\lambda_k)=r_k(\lambda_k)$.}
Differentiating \eqref{eq:factor},
\begin{equation}\label{eq:pprime}
p'(x) \;=\; r_k(x) + (x-\lambda_k)\,r_k'(x).
\end{equation}
Substituting $x=\lambda_k$, the second term vanishes, so $p'(\lambda_k)=r_k(\lambda_k)$. Because the roots are pairwise distinct, $r_k(\lambda_k)=\prod_{j\neq k}(\lambda_k-\lambda_j)\neq 0$, hence $p'(\lambda_k)\neq 0$.

\textbf{Step 2: $p''(\lambda_k)=2\,r_k'(\lambda_k)$.}
Differentiating \eqref{eq:pprime},
\begin{equation}\label{eq:psecond}
p''(x) \;=\; 2\,r_k'(x) + (x-\lambda_k)\,r_k''(x).
\end{equation}
Substituting $x=\lambda_k$, the second term vanishes, giving $p''(\lambda_k)=2\,r_k'(\lambda_k)$.

\textbf{Step 3: A direct computation of $r_k'(\lambda_k)/r_k(\lambda_k)$.}
By the product rule applied to $r_k(x)=\prod_{j\neq k}(x-\lambda_j)$,
\begin{equation}\label{eq:productrule}
r_k'(x) \;=\; \sum_{m\neq k}\prod_{j\neq k,\,j\neq m}(x-\lambda_j).
\end{equation}
For each $m\neq k$, evaluating the corresponding term at $x=\lambda_k$,
\[
\prod_{j\neq k,\,j\neq m}(\lambda_k-\lambda_j) \;=\; \frac{\prod_{j\neq k}(\lambda_k-\lambda_j)}{\lambda_k-\lambda_m} \;=\; \frac{r_k(\lambda_k)}{\lambda_k-\lambda_m},
\]
where the division is permissible because $\lambda_k\neq\lambda_m$. Substituting into \eqref{eq:productrule},
\begin{equation}\label{eq:logd2}
r_k'(\lambda_k) \;=\; \sum_{m\neq k}\frac{r_k(\lambda_k)}{\lambda_k-\lambda_m} \;=\; r_k(\lambda_k)\sum_{m\neq k}\frac{1}{\lambda_k-\lambda_m}.
\end{equation}
Dividing by $r_k(\lambda_k)\neq 0$,
\begin{equation}\label{eq:logderiv}
\frac{r_k'(\lambda_k)}{r_k(\lambda_k)} \;=\; \sum_{j\neq k}\frac{1}{\lambda_k-\lambda_j}.
\end{equation}

\textbf{Step 4: Combine.}
Combining Steps 1, 2, and 3:
\[
\frac{p''(\lambda_k)}{2\,p'(\lambda_k)}
\;=\; \frac{2\,r_k'(\lambda_k)}{2\,r_k(\lambda_k)}
\;=\; \frac{r_k'(\lambda_k)}{r_k(\lambda_k)}
\;=\; \sum_{j\neq k}\frac{1}{\lambda_k-\lambda_j},
\]
which is \eqref{eq:scoreformula}.

\textbf{Step 5: The formula for $\Phi_n(p)$.}
Squaring \eqref{eq:scoreformula} and summing over $k=1,\dots,n$,
\[
\Phi_n(p) \;=\; \sum_{k=1}^n\Bigl(\sum_{j\neq k}\frac{1}{\lambda_k-\lambda_j}\Bigr)^2
\;=\; \sum_{k=1}^n\Bigl(\frac{p''(\lambda_k)}{2\,p'(\lambda_k)}\Bigr)^2
\;=\; \frac{1}{4}\sum_{k=1}^n\Bigl(\frac{p''(\lambda_k)}{p'(\lambda_k)}\Bigr)^2.
\]
This is \eqref{eq:phiformula}.
\end{proof}

\begin{remark}[The case $n=1$]
For $n=1$, every monic degree-$1$ polynomial $p(x)=x-a$ has the single root $a$ and the inner sum $\sum_{j\neq 1}\frac{1}{a-a_j}$ is empty, so $\Phi_1(p)=0$ by the empty-sum convention. With the standard convention $1/0=+\infty$ in $[0,+\infty]$, the inequality \eqref{eq:main} reads $+\infty\geq+\infty+\infty=+\infty$, which holds. The interesting cases are $n\geq 2$.
\end{remark}

\section{Part 2: Full proof for $n=2$ (with equality)}

For $n=2$, the roots are simple as long as $p,q,p\boxplus_2 q$ have nonzero discriminant. We now compute $\Phi_2$ and the convolution explicitly and verify \eqref{eq:main} with equality.

\begin{proposition}\label{prop:n2}
Let $n=2$, let $p(x)=(x-a)(x-b)$ and $q(x)=(x-c)(x-d)$ be monic real-rooted of degree $2$ with $a\neq b$ and $c\neq d$. Then $p\boxplus_2 q$ has discriminant
\begin{equation}\label{eq:disc}
\Delta(p\boxplus_2 q) \;=\; (a-b)^2+(c-d)^2 \;>\; 0,
\end{equation}
so $p\boxplus_2 q$ has two distinct real roots, and
\begin{equation}\label{eq:phi2}
\Phi_2(p) \;=\; \frac{2}{(a-b)^2},\qquad
\Phi_2(q) \;=\; \frac{2}{(c-d)^2},\qquad
\Phi_2(p\boxplus_2 q) \;=\; \frac{2}{(a-b)^2+(c-d)^2}.
\end{equation}
In particular,
\begin{equation}\label{eq:n2eq}
\frac{1}{\Phi_2(p\boxplus_2 q)} \;=\; \frac{1}{\Phi_2(p)}+\frac{1}{\Phi_2(q)},
\end{equation}
so \eqref{eq:main} holds with equality.
\end{proposition}

\begin{proof}
\textbf{Step 1: Coefficients of $p$ and $q$.}
Expanding,
\[
p(x) = x^2 - (a+b)x + ab, \qquad q(x) = x^2 - (c+d)x + cd.
\]
Thus, in the notation of \eqref{eq:ck}, $a_0=b_0=1$, $a_1=-(a+b)$, $a_2=ab$, $b_1=-(c+d)$, $b_2=cd$.

\textbf{Step 2: Compute $c_1$ and $c_2$ from \eqref{eq:ck}.}
Set $n=2$. The relevant ratio in \eqref{eq:ck} is $R(i,j,k):=\frac{(n-i)!(n-j)!}{n!(n-k)!}=\frac{(2-i)!(2-j)!}{2!(2-k)!}$.

For $k=0$: only $(i,j)=(0,0)$, $R=\frac{2!\cdot 2!}{2!\cdot 2!}=1$, so $c_0=a_0b_0=1$ (consistent with monicity).

For $k=1$: pairs $(0,1),(1,0)$.
\begin{itemize}
\item $(0,1)$: $R=\frac{2!\cdot 1!}{2!\cdot 1!}=1$, contribution $a_0b_1=b_1$.
\item $(1,0)$: $R=\frac{1!\cdot 2!}{2!\cdot 1!}=1$, contribution $a_1b_0=a_1$.
\end{itemize}
So
\begin{equation}\label{eq:c1}
c_1 \;=\; a_1+b_1 \;=\; -(a+b+c+d).
\end{equation}

For $k=2$: pairs $(0,2),(1,1),(2,0)$.
\begin{itemize}
\item $(0,2)$: $R=\frac{2!\cdot 0!}{2!\cdot 0!}=1$, contribution $a_0b_2=cd$.
\item $(1,1)$: $R=\frac{1!\cdot 1!}{2!\cdot 0!}=\tfrac{1}{2}$, contribution $\tfrac{1}{2}a_1b_1=\tfrac{1}{2}(a+b)(c+d)$.
\item $(2,0)$: $R=\frac{0!\cdot 2!}{2!\cdot 0!}=1$, contribution $a_2b_0=ab$.
\end{itemize}
So
\begin{equation}\label{eq:c2}
c_2 \;=\; ab+cd+\tfrac{1}{2}(a+b)(c+d).
\end{equation}

\textbf{Step 3: Discriminant of $p\boxplus_2 q$.}
Let $p\boxplus_2 q(x)=x^2+c_1x+c_2$ have roots $\mu_1,\mu_2$. Then $(\mu_1-\mu_2)^2=c_1^2-4c_2$. Compute, using \eqref{eq:c1} and \eqref{eq:c2},
\begin{align}
c_1^2 &= (a+b+c+d)^2 \notag\\
      &= (a+b)^2 + 2(a+b)(c+d) + (c+d)^2,\label{eq:c1sq}\\
4c_2  &= 4ab+4cd+2(a+b)(c+d),\label{eq:4c2}
\end{align}
so
\begin{align}
c_1^2 - 4c_2
&= (a+b)^2 - 4ab + (c+d)^2 - 4cd \notag\\
&= (a-b)^2 + (c-d)^2.\label{eq:discproof}
\end{align}
This proves \eqref{eq:disc}. Since by hypothesis $a\neq b$ and $c\neq d$, both $(a-b)^2$ and $(c-d)^2$ are strictly positive, hence $c_1^2-4c_2>0$. Therefore $\mu_1,\mu_2\in\mathbb{R}$ are distinct.

\textbf{Step 4: Evaluate $\Phi_2$.}
For a degree-$2$ monic polynomial $r(x)=(x-\mu_1)(x-\mu_2)$ with $\mu_1\neq\mu_2$, the score formula \eqref{eq:scoreformula} of Lemma~\ref{lem:score} gives
\[
\sum_{j\neq 1}\frac{1}{\mu_1-\mu_j} = \frac{1}{\mu_1-\mu_2},\qquad
\sum_{j\neq 2}\frac{1}{\mu_2-\mu_j} = \frac{1}{\mu_2-\mu_1} = -\frac{1}{\mu_1-\mu_2}.
\]
Hence
\begin{equation}\label{eq:phi2gen}
\Phi_2(r) \;=\; \Bigl(\frac{1}{\mu_1-\mu_2}\Bigr)^2 + \Bigl(\frac{1}{\mu_2-\mu_1}\Bigr)^2 \;=\; \frac{2}{(\mu_1-\mu_2)^2}.
\end{equation}
Applying \eqref{eq:phi2gen} to $p$, $q$, and $p\boxplus_2 q$ together with $(\mu_1-\mu_2)^2=(a-b)^2+(c-d)^2$ from Step~3, we obtain \eqref{eq:phi2}.

\textbf{Step 5: Equality \eqref{eq:n2eq}.}
From \eqref{eq:phi2},
\[
\frac{1}{\Phi_2(p)}+\frac{1}{\Phi_2(q)}
= \frac{(a-b)^2}{2} + \frac{(c-d)^2}{2}
= \frac{(a-b)^2+(c-d)^2}{2}
= \frac{1}{\Phi_2(p\boxplus_2 q)}.
\]
This is exactly \eqref{eq:n2eq}.
\end{proof}

\begin{remark}[Boundary cases for $n=2$]
If $a=b$ (so $p$ has a double root), then $\Phi_2(p)=+\infty$ by definition, and $1/\Phi_2(p)=0$. The same logic as Step~4 of Proposition~\ref{prop:n2} still gives $\Phi_2(p\boxplus_2 q)=2/((a-b)^2+(c-d)^2)=2/(c-d)^2$ provided $c\neq d$, and the inequality becomes $1/\Phi_2(p\boxplus_2 q)=(c-d)^2/2 = 0+1/\Phi_2(q)$, so equality still holds. If both $a=b$ and $c=d$, both sides are $+\infty\geq+\infty+\infty$, interpreted in $[0,+\infty]$ as $0\geq 0+0$, which holds. Hence \eqref{eq:main} holds with equality on the entire space of pairs of degree-$2$ monic real-rooted polynomials.
\end{remark}

\section{Part 3: The general case $n\geq 3$ --- partial result and identification of the missing ingredient}

The inequality \eqref{eq:main} for general $n$ is the finite-free analogue of a deep theorem in free probability: Voiculescu's superadditivity of the inverse free Fisher information under free additive convolution,
\begin{equation}\label{eq:voic}
\frac{1}{\Phi^*(\mu\boxplus\nu)} \;\geq\; \frac{1}{\Phi^*(\mu)}+\frac{1}{\Phi^*(\nu)},
\end{equation}
proved by Voiculescu (1998) for free Fisher information $\Phi^*$ of probability measures, see also Hiai--Petz, \emph{The Semicircle Law, Free Random Variables and Entropy} (AMS, 2000), Ch.~7. The connection comes from the fact that, under appropriate scaling, the empirical root distribution of $p\boxplus_n q$ converges to the free additive convolution of those of $p$ and $q$ as $n\to\infty$, and the per-particle score $\sum_{j\neq i}\frac{1}{\lambda_i-\lambda_j}$ approximates the Hilbert transform of the empirical density used in the definition of $\Phi^*$. Even the asymptotic statement is nontrivial; the finite-$n$ statement \eqref{eq:main} is more delicate and, to our knowledge, requires an argument tailored to finite free probability rather than a direct quotation of \eqref{eq:voic}.

We now (i) record an inequality that follows directly from our Part~2 computation and (ii) state precisely what the missing ingredient for $n\geq 3$ is.

\subsection{Auxiliary observations and the strongest unconditional partial result}

Let $\epsilon>0$ and consider the operator $T_\epsilon=(1-\epsilon\,d/dx)^n$ from the problem statement. Acting on a monic polynomial $p$ of degree $n$, $T_\epsilon p$ is again a monic polynomial of degree $n$: writing $p(x)=\sum_{k=0}^n a_k x^{n-k}$,
\[
T_\epsilon p(x) \;=\; \sum_{m=0}^n\binom{n}{m}(-\epsilon)^m p^{(m)}(x),
\]
and the coefficient of $x^n$ is $a_0=1$. The role of $T_\epsilon$ in proving \eqref{eq:main} for $n\geq 3$ is discussed under ingredient (III) below.

The next claim is the strongest unconditional partial result we are able to verify with full proof.

\begin{proposition}[Cauchy--Schwarz lower bound on $\Phi_n$]\label{prop:cs}
Let $p$ be a monic real-rooted polynomial of degree $n\geq 2$ with simple roots $\lambda_1,\dots,\lambda_n$. Then
\begin{equation}\label{eq:cs}
\Phi_n(p) \;\geq\; \frac{1}{n}\Bigl(\sum_{i=1}^n\sum_{j\neq i}\frac{1}{\lambda_i-\lambda_j}\Bigr)^2 \;=\; 0.
\end{equation}
\end{proposition}
\begin{proof}
By Cauchy--Schwarz,
\[
\Phi_n(p) = \sum_{i=1}^n s_i^2 \;\geq\; \frac{1}{n}\Bigl(\sum_{i=1}^n s_i\Bigr)^2,\qquad s_i:=\sum_{j\neq i}\frac{1}{\lambda_i-\lambda_j}.
\]
The double sum $\sum_i\sum_{j\neq i}\frac{1}{\lambda_i-\lambda_j}$ pairs $(i,j)$ with $(j,i)$, contributing $\frac{1}{\lambda_i-\lambda_j}+\frac{1}{\lambda_j-\lambda_i}=0$. Hence the right side of \eqref{eq:cs} is $0$, recovering the trivial bound $\Phi_n(p)\geq 0$.
\end{proof}

The bound \eqref{eq:cs} is too weak to prove \eqref{eq:main}; we record it only to delineate exactly what we have established unconditionally and what we have not.

\subsection{Identification of the missing ingredient for $n\geq 3$}

To extend the equality of Proposition~\ref{prop:n2} to an inequality for $n\geq 3$ requires, at minimum, one of the following three new ingredients, none of which we are able to supply with a complete proof here:

\begin{enumerate}
\item[(I)] \textbf{Convexity of $1/\Phi_n$ along finite-free convolution paths.}
\\Define, for $t\in[0,1]$, the path $r_t := p\boxplus_n q_t$ where $q_t$ interpolates $q$ to a degenerate (Dirac-type) limit. If
\[
t \;\mapsto\; \frac{1}{\Phi_n(r_t)}
\]
is convex on $[0,1]$, then a standard chain-rule argument (cf.~\cite[Prop.~7.7]{HiaiPetz} in the continuous setting) gives \eqref{eq:main}. We have not been able to verify this convexity in finite-free probability without additional input.

\item[(II)] \textbf{A finite-free chain rule for relative Fisher information.}
\\In the continuous free setting, $\Phi^*(X+Y)$ is controlled by $\Phi^*(X)$ and $\Phi^*(Y)$ through a chain rule for conditional expectation onto the algebra generated by $X+Y$. The analogue in finite free probability would be an identity of the form
\[
s_i(p\boxplus_n q) \;=\; \mathbb{E}\bigl[s_{\sigma(i)}(p) + s_{\tau(i)}(q)\,\big|\,\mathcal{F}_{p\boxplus_n q}\bigr]
\]
for some appropriate finite-dimensional conditional expectation. We do not have such an identity; we are not aware of it being established in the literature.

\item[(III)] \textbf{Monotonicity of $\Phi_n$ under $T_\epsilon=(1-\epsilon\,d/dx)^n$.}
\\If one can prove $\Phi_n(T_\epsilon p)\leq \Phi_n(p)$ together with a one-parameter semigroup structure of $T_\epsilon$ that intertwines $\boxplus_n$, then a Bakry--\'Emery-type entropy production argument can convert the equality \eqref{eq:n2eq} into the inequality \eqref{eq:main} for all $n$. The required intertwining and the monotonicity statement appear plausible (the analogous statements for the free Ornstein--Uhlenbeck semigroup are classical), but we cannot establish them without a substantive new computation.
\end{enumerate}

In all three cases, the gap is not a routine extension of Part~2. The $n=2$ proof exploits the fact that the discriminant of a quadratic is a quadratic form in the coefficients with the special property
$
c_1^2-4c_2=(a-b)^2+(c-d)^2,
$
which is exactly an additive Pythagorean identity. For $n\geq 3$ the discriminant of $p\boxplus_n q$ becomes a polynomial of degree $2(n-1)$ in the coefficients, with no analogous decomposition into a sum of squares of root differences of $p$ and of $q$ separately. A direct generalization of the algebraic proof in Part~2 therefore fails.

\subsection{Status of \eqref{eq:main} for $n\geq 3$}

In summary, we have proved \eqref{eq:main} in full for $n=1$ and $n=2$, with equality. For $n\geq 3$ we have established a careful reduction (Lemma~\ref{lem:score}) and identified three precise mathematical ingredients (I), (II), (III) any one of which would suffice. The full statement \eqref{eq:main} for $n\geq 3$ is, to the best of our knowledge, equivalent to the finite-free analogue of Voiculescu's free Fisher information superadditivity \eqref{eq:voic}, which is a deep theorem requiring tools from free stochastic analysis whose finite-free counterparts are not yet developed in the literature available to us.

\section*{Conclusion}

\begin{itemize}
\item For $n=1$ and $n=2$ the inequality \eqref{eq:main} holds with equality (Proposition~\ref{prop:n2}); the answer to the binary question for these cases is \textbf{Yes}.
\item For $n\geq 3$ we conjecture, but cannot prove with full rigor here, that \eqref{eq:main} continues to hold; we have identified three precise auxiliary statements (convexity of $1/\Phi_n$ along $\boxplus_n$ paths; a finite-free chain rule for relative Fisher information; monotonicity of $\Phi_n$ under $T_\epsilon$) any one of which would suffice. We do not claim a full proof for $n\geq 3$.
\end{itemize}

\begin{thebibliography}{9}
\bibitem{Voiculescu}
D.~Voiculescu, \emph{The analogues of entropy and of Fisher's information measure in free probability theory, V: Noncommutative Hilbert transforms}, Invent.\ Math.\ \textbf{132} (1998), 189--227.

\bibitem{HiaiPetz}
F.~Hiai and D.~Petz, \emph{The Semicircle Law, Free Random Variables and Entropy}, Mathematical Surveys and Monographs, vol.~77, AMS, Providence, RI, 2000.

\bibitem{MSS}
A.~Marcus, D.~Spielman, and N.~Srivastava, \emph{Finite free convolutions of polynomials}, Probab.\ Theory Related Fields \textbf{182} (2022), 807--848.
\end{thebibliography}

\end{document}
